\documentclass[11pt,a4paper]{article}
\usepackage[T1]{fontenc}
\usepackage[utf8]{inputenc}
\usepackage{amsmath,amssymb,amsthm}
\usepackage[margin=1in]{geometry}
\usepackage[colorlinks=true,linkcolor=blue,urlcolor=blue]{hyperref}
\theoremstyle{plain}
\newtheorem{theorem}{Theorem}
\newtheorem{lemma}[theorem]{Lemma}
\newtheorem{proposition}[theorem]{Proposition}
\theoremstyle{definition}
\newtheorem{remark}[theorem]{Remark}
\title{A proof of the conjectured generating function for OEIS A294139}
\author{Adrian Perez Fontelles\\ \small Independent researcher}
\date{31 August 2026}
\begin{document}
\maketitle

\begin{abstract}
OEIS A294139 carries a conjectured generating function. We prove it. The entry states, as fact,
a closed form, and from that a linear recurrence with polynomial coefficients satisfied by
the sequence is derived exactly. The conjectured generating function is then shown to satisfy the
same recurrence, and to agree with the sequence at enough consecutive indices. A recurrence
of order $4$ whose leading coefficient does not vanish, together with $4$
consecutive values, determines a sequence completely, so the two agree everywhere. Nothing
is checked numerically and then assumed: the recurrence is derived symbolically, the
verification is an exact identity, and the finitely many indices where the leading
coefficient vanishes are computed rather than waved at.
\end{abstract}

\noindent\small 2020 Mathematics Subject Classification. 05A15, 11B37, 33F10.\normalsize

\section{The sequence, what is known, and what is conjectured}

OEIS A294139 is ``Sum of the areas of the distinct rectangles (and the areas of the squares on their sides) with positive integer sides such that L + W = n, W < L''. It has offset $1$ and begins
\[
a(1),\dots,a(8)\;=\;0, 0, 12, 23, 70, 105, 210, 282,\ \dots
\]
The entry states, not as a conjecture,
\begin{quote}\small
a(n) = n*(4-21*n+12*n\textasciicircum{}2-5*n*(-1)\textasciicircum{}n)/16. - \emph{Wesley Ivan Hurt}, Dec 02 2023
\end{quote}
and separately carries the comment
\begin{quote}\small
G.f.: x\textasciicircum{}3*(12 + 11*x + 11*x\textasciicircum{}2 + 2*x\textasciicircum{}3) / ((1 - x)\textasciicircum{}4*(1 + x)\textasciicircum{}3).
\end{quote}
As of the ``Last modified'' line on the live entry (2024-10-14, revision 49) the second
statement is still recorded as a conjecture, and no proof appears among the entry's links,
formulas or programs.

\section{A recurrence from what is known}

The closed form the entry posts is a sum of finitely many hypergeometric terms. For terms $t_{1},\dots,t_{m}$ one may seek $c_{0},\dots,c_{m}$ in $\mathbb{Q}(n)$ with $\sum_{k}c_{k}(n)t_{j}(n+k)=0$ for every $j$; dividing the $j$-th equation by $t_{j}(n)$ makes every coefficient a rational function, so this is $m$ equations in $m+1$ unknowns over $\mathbb{Q}(n)$ and a nonzero solution exists. It gives

\begin{equation}\label{eq:rec}
\left(- 2 n^{5} - 25 n^{4} - 120 n^{3} - 275 n^{2} - 298 n - 120\right)\,a(n) + \left(4 n^{5} + 48 n^{4} + 212 n^{3} + 408 n^{2} + 288 n\right)\,a(n+1) + \left(- 6 n^{4} - 48 n^{3} - 114 n^{2} - 72 n\right)\,a(n+2) + \left(- 4 n^{5} - 32 n^{4} - 84 n^{3} - 88 n^{2} - 32 n\right)\,a(n+3) + \left(2 n^{5} + 15 n^{4} + 40 n^{3} + 45 n^{2} + 18 n\right)\,a(n+4) \;=\; 0
\end{equation}

The leading coefficient of \eqref{eq:rec} is $n \left(n + 1\right) \left(n + 2\right) \left(n + 3\right) \left(2 n + 3\right)$. Its only integer zeros are $n\in\{-3, -2, -1, 0\}$, all below the range claimed here.

\section{The conjectured description satisfies it}

Write $\theta=x\,\frac{d}{dx}$ and, for the coefficients $p_{i}$ of \eqref{eq:rec} written in the backward form $\sum_{i}p_{i}(n)a(n-i)=0$, set $B(x)=\sum_{i}x^{i}\,(p_{i}(\theta+i)A)(x)$ for the conjectured generating function $A$. Then $[x^{n}]B=\sum_{i}p_{i}(n)a(n-i)$, so the conjectured generating function has coefficients satisfying \eqref{eq:rec} exactly when $B$ is a polynomial. Reducing $B$ inside the field $A$ generates, which is closed under $\theta$, gives a polynomial, and the verification is an exact cancellation of rational functions.

\section{Uniqueness}

\begin{lemma}\label{lem:unique}
Let $\sum_{j=0}^{r}q_{j}(n)\,b(n+j)=0$ be a linear recurrence whose leading coefficient
$q_{r}$ has no zero at any integer $n\ge n_{0}$. If two sequences both satisfy it for all
$n\ge n_{0}$ and agree at the $r$ consecutive indices $n_{0},\dots,n_{0}+r-1$, they agree
at every $n\ge n_{0}$.
\end{lemma}

\begin{proof}
Induction. Suppose the two agree at $n_{0},\dots,m+r-1$ for some $m\ge n_{0}$. The
recurrence at $n=m$ determines $b(m+r)$ from the preceding $r$ values, because $q_{r}(m)$
is nonzero and may be divided by; both sequences therefore take the same value at $m+r$.
\end{proof}

\begin{theorem}
For all $n\ge 1$,
\[
\sum_{n\ge 1} a(n)x^{n} \;=\; \frac{x^{3} \left(2 x^{3} + 11 x^{2} + 11 x + 12\right)}{\left(x - 1\right)^{4} \left(x + 1\right)^{3}}
\]
\end{theorem}

\begin{proof}
Both sides satisfy \eqref{eq:rec}: the sequence because the recurrence was derived from
the entry's own a closed form, and the conjectured generating function by the computation of the
previous section. They agree at the $4$ consecutive indices beginning
$n=1$, which was checked against the entry's DATA in exact integer arithmetic.
The leading coefficient has no zero at any integer $n\ge1$, so
Lemma~\ref{lem:unique} applies.
\end{proof}

\section{Verification}

Three checks, independent of one another.

First, the description the entry states as fact was itself confirmed against the entry's
DATA before any recurrence was derived from it; a statement that did not reproduce the
entry's own terms would not have been used as the basis of anything.

Second, the derived recurrence \eqref{eq:rec} was evaluated on the published terms in
exact integer arithmetic and vanishes at every index where it applies.

Third, the conjectured generating function was evaluated at the first $13$ indices and
compared term by term with the DATA; the values agree exactly. This is not part of the
proof, which is symbolic, but it would catch a transcription error in either statement.

\begin{thebibliography}{9}
\bibitem{oeis} The OEIS Foundation, \emph{The On-Line Encyclopedia of Integer Sequences},
\url{https://oeis.org}, 2026. Sequence A294139.
\bibitem{aeqb} M.~Petkov\v{s}ek, H.~S.~Wilf and D.~Zeilberger, \emph{A=B}, A K Peters,
1996.
\bibitem{fs} P.~Flajolet and R.~Sedgewick, \emph{Analytic Combinatorics}, Cambridge
University Press, 2009.
\bibitem{kauers} M.~Kauers and P.~Paule, \emph{The Concrete Tetrahedron}, Springer, 2011.
\end{thebibliography}

\end{document}
